% GENERATED FILE. Edit canonical lessons, then run npm run generate:books.
% canonical-source-hash: fnv1a64:419e926fbac3c101
% canonical-lessons: RU-C453-primerochnaya, RU-C453-primeryat, RU-C453-razmer, RU-C453-zhat, RU-C453-obmenivat

\chapter{Trying things on}
\label{ch:ru-trying-things-on}
\hlchaptermodality{453}

\begin{chapteropening}
\textbf{By the end of this chapter:} \emph{I can try on clothes: the fitting room, trying on, the size, shoes that pinch and exchanging things.}

Complete the last lesson of chapter 453: I can try on clothes: the fitting room, trying on, the size, shoes that pinch and exchanging things.
\end{chapteropening}

\section[\texorpdfstring{\ru{примерочная} (primérochnaya)}{примерочная (primérochnaya)}]{\ru{примерочная} (primérochnaya) — a fitting room}

\label{lesson:RU-C453-primerochnaya}



\begin{quote}
Before the new one: say the Russian for a shop assistant, then the Russian for to weigh.
\end{quote}

\subsection*{You'll want to know: \ru{примерочная}}
\textbf{\ru{примерочная}} (\emph{primérochnaya}) — \textquotedblleft{}a fitting room\textquotedblright{}.

\begin{quote}
\ru{Примерочная} \ru{слева} \ru{от} \ru{кассы}.
\end{quote}

\begin{quote}
\emph{Primérochnaya sléva ot kássy.}
\end{quote}

\begin{quote}
The fitting room is to the left of the till.
\end{quote}

An adjective used as a noun.

\begin{cousinweb}
\emph{\ru{примерочная}} is built on \emph{\ru{примерка}} 'a fitting', from \emph{\ru{примерять}} 'to try on', on \emph{\ru{мера}} 'a measure'. EtymDB builds \emph{\ru{мерить}} on Proto-Slavic \emph{měra}.
\end{cousinweb}

\begin{grammarlens}[title={What you've built}]
One more word for telling, explaining and describing at B1.
\end{grammarlens}

\subsection*{Your turn}
\begin{itemize}
\raggedright
  \item \emph{Say it:} \emph{primérochnaya}
  \item \emph{Say it:} \emph{primérochnaya}, once more
  \item \emph{Say it:} a sentence of your own with \emph{primérochnaya}
  \item \emph{Recall:} say the Russian for a shop assistant, then the Russian for to weigh, then say \emph{primérochnaya} again
\end{itemize}

\begin{tcolorbox}[breakable,colback=teal!4,colframe=teal!35!black,title={Before you move on}]
What does \ru{примерочная} mean? (\textquotedblleft{}A fitting room\textquotedblright{}.) Say it once more.
\end{tcolorbox}
\section[\texorpdfstring{\ru{примерять} (primeryát)}{примерять (primeryát)}]{\ru{примерять} (primeryát) — to try on}

\label{lesson:RU-C453-primeryat}



\begin{quote}
Before the new one: say the Russian for to weigh, then the Russian for a fitting room.
\end{quote}

\subsection*{You'll want to know: \ru{примерять}}
\textbf{\ru{примерять}} (\emph{primeryát}) — \textquotedblleft{}to try on\textquotedblright{}.

\begin{quote}
\ru{Она} \ru{долго} \ru{примеряла} \ru{платья}.
\end{quote}

\begin{quote}
\emph{Oná dólgo primeryála plátya.}
\end{quote}

\begin{quote}
She spent a long time trying on dresses.
\end{quote}

Imperfective; perfective \emph{\ru{примерить}} (\emph{\ru{примерю}, \ru{примеришь}}).

\begin{cousinweb}
\emph{\ru{примерять}} is built from \emph{\ru{при}-} and \emph{\ru{мерять}}, \emph{\ru{мерить}} 'to measure', which EtymDB builds on Proto-Slavic \emph{měra} 'a measure'.
\end{cousinweb}

\begin{grammarlens}[title={What you've built}]
One more word for telling, explaining and describing at B1.
\end{grammarlens}

\subsection*{Your turn}
\begin{itemize}
\raggedright
  \item \emph{Say it:} \emph{primeryát}
  \item \emph{Say it:} \emph{primeryát}, once more
  \item \emph{Say it:} a sentence of your own with \emph{primeryát}
  \item \emph{Recall:} say the Russian for to weigh, then the Russian for a fitting room, then say \emph{primeryát} again
\end{itemize}

\begin{tcolorbox}[breakable,colback=teal!4,colframe=teal!35!black,title={Before you move on}]
What does \ru{примерять} mean? (\textquotedblleft{}To try on\textquotedblright{}.) Say it once more.
\end{tcolorbox}
\section[\texorpdfstring{\ru{размер} (razmér)}{размер (razmér)}]{\ru{размер} (razmér) — a size}

\label{lesson:RU-C453-razmer}



\begin{quote}
Before the new one: say the Russian for a fitting room, then the Russian for to try on.
\end{quote}

\subsection*{You'll want to know: \ru{размер}}
\textbf{\ru{размер}} (\emph{razmér}) — \textquotedblleft{}a size\textquotedblright{}.

\begin{quote}
\ru{Эти} \ru{брюки} \ru{мне} \ru{не} \ru{подошли} \ru{по} \ru{размеру}.
\end{quote}

\begin{quote}
\emph{Éti bryúki mne ne podoshlí po razméru.}
\end{quote}

\begin{quote}
These trousers were not my size.
\end{quote}

Masculine. \emph{\ru{Какой} \ru{у} \ru{вас} \ru{размер}?} — what size are you?

\begin{cousinweb}
\emph{\ru{размер}} is built from \emph{\ru{раз}-} and \emph{\ru{мера}} 'a measure'. EtymDB lists \emph{\ru{размер}} without a deeper source.
\end{cousinweb}

\begin{grammarlens}[title={What you've built}]
One more word for telling, explaining and describing at B1.
\end{grammarlens}

\subsection*{Your turn}
\begin{itemize}
\raggedright
  \item \emph{Say it:} \emph{razmér}
  \item \emph{Say it:} \emph{razmér}, once more
  \item \emph{Say it:} a sentence of your own with \emph{razmér}
  \item \emph{Recall:} say the Russian for a fitting room, then the Russian for to try on, then say \emph{razmér} again
\end{itemize}

\begin{tcolorbox}[breakable,colback=teal!4,colframe=teal!35!black,title={Before you move on}]
What does \ru{размер} mean? (\textquotedblleft{}A size\textquotedblright{}.) Say it once more.
\end{tcolorbox}
\section[\texorpdfstring{\ru{жать} (zhat)}{жать (zhat)}]{\ru{жать} (zhat) — to pinch, to be too tight; to press}

\label{lesson:RU-C453-zhat}



\begin{quote}
Before the new one: say the Russian for to try on, then the Russian for a size.
\end{quote}

\subsection*{You'll want to know: \ru{жать}}
\textbf{\ru{жать}} (\emph{zhat}) — \textquotedblleft{}to pinch, to be too tight; to press\textquotedblright{}.

\begin{quote}
\ru{Новые} \ru{туфли} \ru{немного} \ru{жмут}.
\end{quote}

\begin{quote}
\emph{Nóvyye túfli nemnógo zhmut.}
\end{quote}

\begin{quote}
The new shoes pinch a little.
\end{quote}

Imperfective. Present \emph{\ru{жму}, \ru{жмёшь}, \ru{жмут}}. \emph{\ru{Жать} \ru{руку}} — to shake hands. A different verb \emph{\ru{жать}} (\emph{\ru{жну}}) means to reap.

\begin{cousinweb}
\emph{\ru{жать}} 'to press, to squeeze' is usually counted an inherited Slavic verb. EtymDB lists it without a source, so its deeper origin is left open here.
\end{cousinweb}

\begin{grammarlens}[title={What you've built}]
One more word for telling, explaining and describing at B1.
\end{grammarlens}

\subsection*{Your turn}
\begin{itemize}
\raggedright
  \item \emph{Say it:} \emph{zhat}
  \item \emph{Say it:} \emph{zhat}, once more
  \item \emph{Say it:} a sentence of your own with \emph{zhat}
  \item \emph{Recall:} say the Russian for to try on, then the Russian for a size, then say \emph{zhat} again
\end{itemize}

\begin{tcolorbox}[breakable,colback=teal!4,colframe=teal!35!black,title={Before you move on}]
What does \ru{жать} mean? (\textquotedblleft{}To pinch, to be too tight; to press\textquotedblright{}.) Say it once more.
\end{tcolorbox}
\section[\texorpdfstring{\ru{обменивать} (obménivat)}{обменивать (obménivat)}]{\ru{обменивать} (obménivat) — to exchange (goods), to swap}

\label{lesson:RU-C453-obmenivat}



\begin{quote}
Before the new one: say the Russian for a size, then the Russian for to pinch.
\end{quote}

\subsection*{You'll want to know: \ru{обменивать}}
\textbf{\ru{обменивать}} (\emph{obménivat}) — \textquotedblleft{}to exchange (goods), to swap\textquotedblright{}.

\begin{quote}
\ru{В} \ru{этом} \ru{магазине} \ru{легко} \ru{обменивают} \ru{товар}.
\end{quote}

\begin{quote}
\emph{V étom magazíne legkó obménivayut továr.}
\end{quote}

\begin{quote}
In this shop they exchange goods without any fuss.
\end{quote}

Imperfective; perfective \emph{\ru{обменять}}. Takes \emph{\ru{на}} + accusative: \emph{\ru{обменять} \ru{свитер} \ru{на} \ru{другой} \ru{размер}}.

\begin{cousinweb}
\emph{\ru{обменивать}} is built from \emph{\ru{об}-} and \emph{-\ru{менивать}}, from \emph{\ru{менять}} 'to change', on \emph{\ru{мена}} 'an exchange', which is inherited from Proto-Slavic \emph{měna} (EtymDB).
\end{cousinweb}

\begin{grammarlens}[title={What you've built}]
One more word for telling, explaining and describing at B1.
\end{grammarlens}

\subsection*{Your turn}
\begin{itemize}
\raggedright
  \item \emph{Say it:} \emph{obménivat}
  \item \emph{Say it:} \emph{obménivat}, once more
  \item \emph{Say it:} a sentence of your own with \emph{obménivat}
  \item \emph{Recall:} say the Russian for a size, then the Russian for to pinch, then say \emph{obménivat} again
\end{itemize}

\begin{tcolorbox}[breakable,colback=teal!4,colframe=teal!35!black,title={Before you move on}]
What does \ru{обменивать} mean? (\textquotedblleft{}To exchange (goods), to swap\textquotedblright{}.) Say it once more.
\end{tcolorbox}
